\documentclass[11pt,a4paper]{article}
% ==========================================================
% T0_preamble_local_En.tex
% T0/FFGFT Documents — Kindle 6×9 Format (English)
% Basiert auf harmonikabau_preamble.sty
% Schrift: DejaVu Sans | Compiler: xelatex
% ==========================================================

% --- Schrift ---
\RequirePackage{fontspec}
\setmainfont{DejaVu Sans}
\setsansfont{DejaVu Sans}
\setmonofont{DejaVu Sans Mono}[Scale=0.88]
\defaultfontfeatures{Ligatures=TeX}

% --- Sprache ---
\RequirePackage{polyglossia}
\setmainlanguage{english}

% --- Mathematik ---
\RequirePackage{amsmath,amssymb,amsthm}
\RequirePackage{mathtools}
\RequirePackage{unicode-math}
\setmathfont{Latin Modern Math}[Scale=1.0]
\RequirePackage{siunitx}
\sisetup{locale=UK, per-mode=fraction, separate-uncertainty=true}

% --- Seitenlayout: direkt 6×9 Zoll mit KDP-konformen Rändern ---
\RequirePackage{geometry}
\geometry{
  paperwidth  = 6in,
  paperheight = 9in,
  left        = 19.1mm,
  right       = 12.7mm,
  top         = 14mm,
  bottom      = 14mm,
  headheight  = 10pt,
  headsep     = 3mm,
  footskip    = 4mm
}
% Verhindert sub-pixel Rundungsfehler die KDP als Randverstoß erkennt
\addtolength{\textwidth}{-2pt}

% --- Tabellen ---
\RequirePackage{booktabs}
\RequirePackage{colortbl}
\RequirePackage{array}
\RequirePackage{tabularx}
\RequirePackage{longtable}
\RequirePackage{multirow}

% --- Grafik & Farbe ---
\RequirePackage{graphicx}
\RequirePackage{xcolor}
\RequirePackage{float}
\RequirePackage{caption}
\captionsetup{font=small, labelfont=bf, labelsep=period}

% --- Listen ---
\RequirePackage{enumitem}
\setlist[itemize]{topsep=2pt, itemsep=1pt, parsep=0pt}
\setlist[enumerate]{topsep=2pt, itemsep=1pt, parsep=0pt}

% --- Boxen ---
\RequirePackage[many]{tcolorbox}

% --- Links ---
\RequirePackage{hyperref}
\hypersetup{colorlinks=true, linkcolor=T0blue, citecolor=T0green, urlcolor=T0blue!70!black}

% --- Zeilenumbruch: sauber wie im Harmonikabuch ---
\sloppy
\emergencystretch=3em
\tolerance=2500
\hbadness=10000
\hfuzz=2pt
\clubpenalty=10000
\widowpenalty=10000

% ============================================================
% Farben
% ============================================================
\definecolor{T0blue}{HTML}{16213e}
\definecolor{T0red}{HTML}{e94560}
\definecolor{T0green}{HTML}{1b5e20}
\definecolor{T0orange}{HTML}{e65100}
\definecolor{T0gray}{HTML}{616161}
\definecolor{T0lightgray}{HTML}{f5f5f5}
\definecolor{T0headbg}{HTML}{dce8f5}

% ============================================================
% Überschriften (angepasst auf 6×9)
% ============================================================
\RequirePackage{titlesec}

\titleformat{\section}
  {\large\bfseries\color{T0blue}}
  {\thesection}{1em}{}
\titleformat{\subsection}
  {\normalsize\bfseries\color{T0blue!80}}
  {\thesubsection}{1em}{}
\titleformat{\subsubsection}
  {\small\bfseries\color{T0gray}}
  {\thesubsubsection}{1em}{}

\titlespacing*{\section}{0pt}{10pt}{4pt}
\titlespacing*{\subsection}{0pt}{7pt}{2pt}
\titlespacing*{\subsubsection}{0pt}{5pt}{1pt}

% ============================================================
% Kopf- und Fußzeile
% ============================================================
\RequirePackage{fancyhdr}
\pagestyle{fancy}
\fancyhf{}
\renewcommand{\headrulewidth}{0.4pt}
\renewcommand{\footrulewidth}{0pt}
\fancyhead[L]{\small\color{T0gray}\docheadleft}
\fancyhead[R]{\small\color{T0gray}\thepage}
\fancyfoot[C]{\small\color{T0gray}\docfootcenter}

\fancypagestyle{firstpage}{
  \fancyhf{}
  \renewcommand{\headrulewidth}{0pt}
  \fancyfoot[C]{\small\color{T0gray}\docfootcenter}
}

% ============================================================
% Boxen (kompatibel mit ch-Dateien)
% ============================================================
\newtcolorbox{keybox}[1][Key Point]{%
  colback=T0green!7, colframe=T0green!55!black,
  boxrule=0.5pt, arc=2pt,
  left=3mm, right=3mm, top=1.5mm, bottom=1.5mm,
  fontupper=\small, title={#1},}

\newtcolorbox{warnbox}[1][Note]{%
  colback=T0red!7, colframe=T0red!60!black,
  boxrule=0.5pt, arc=2pt,
  left=3mm, right=3mm, top=1.5mm, bottom=1.5mm,
  fontupper=\small, title={#1},}

\newtcolorbox{infobox}[1][Information]{%
  colback=T0lightgray, colframe=T0gray!40,
  boxrule=0.4pt, arc=2pt,
  left=3mm, right=3mm, top=1.5mm, bottom=1.5mm,
  fontupper=\small, title={#1},}

\newtcolorbox{predbox}[1][Prediction]{%
  colback=T0blue!7, colframe=T0blue!50,
  boxrule=0.5pt, arc=2pt,
  left=3mm, right=3mm, top=1.5mm, bottom=1.5mm,
  fontupper=\small, title={#1},}

\newtcolorbox{important}[1][Important]{%
  colback=T0orange!7, colframe=T0orange!60!black,
  boxrule=0.5pt, arc=2pt,
  left=3mm, right=3mm, top=1.5mm, bottom=1.5mm,
  fontupper=\small, title={#1},}

% ============================================================
% Metadaten-Befehle (werden im Wrapper gesetzt)
% ============================================================
\newcommand{\docheadleft}{}
\newcommand{\docfootcenter}{}

% Zeilenabstand
\RequirePackage{setspace}
\setstretch{1.12}

\hfuzz=8pt\tolerance=5000\emergencystretch=3em
\title{Algebraic Bridges to the Consciousness Discussion\\
\large Connection to Documents 100 and 170}
\author{Johann Pascher\\
\small Independent Researcher, Upper Austria, Austria\\
\small ORCID: \href{https://orcid.org/0009-0000-6518-4064}{0009-0000-6518-4064}\\
\small \href{mailto:johann.pascher@gmail.com}{johann.pascher@gmail.com}}
\date{8 May 2026}
\begin{document}
\maketitle
\begin{abstract}
\noindent
Document 206 (\emph{Triangle--Matrix Reduction Theorem}) provides an algebraic structure that grounds the existing consciousness discussion in the FFGFT series geometrically, without extending it in content. The phenomenological line is already drawn in Document~100 (\emph{Consciousness as fractal incoherence}, connecting to the no-go theorem of Adlam, McQueen and Waegell, 2025) and in Document~170 (\emph{Consciousness as a topologically forced threshold ladder}).

We show three structural bridges: (A) Tekermen's \emph{constitutive openness} is algebraically identical to the non-closure theorem (Doc~193) and corresponds to the condition $\xi > 0$; a perfectly unitary system has $\det A = 1$, hence no world-relation. (B) Phillips' \emph{Information Grounding Law} ($A^2 = A^\top$) is the Z$_3$ self-inversion and corresponds to the \emph{recursive sensory feedback} of Document~170. (C) The Z$_3^3$ scale ladder with eigenvalues $(1-\xi)^{n/3}$ provides the algebraic skeleton for the discrete complexity steps that Document~170 formulates as a working hypothesis; it enforces discreteness, but not the phenomenological step assignment.

We sharply delimit what the algebra does \emph{not} achieve: no derivation of consciousness from $\xi$, no prediction of which layer $n/3$ is conscious, no resolution of the Hard~Problem. The bridges are structural, not ontological. A methodological section transfers the self-checks from Document~206 (numerology filter, circularity section) to the consciousness topic. An additional chapter formulates four qualitative empirical predictions under the discreteness hypothesis from Document~170 --- bimodality, developmental jump, anaesthesia tipping, species classification --- each without quantitative values derived from $\xi$.

The balance: Document~206 makes the language of the discussion algebraically consistent, but the consciousness question itself remains an open research question; it is neither solved nor closed by Doc~206.
\end{abstract}
\tableofcontents
\newpage
% ============================================================
% dok207_bewusstsein_bruecken_En_ch.tex
% Algebraic Bridges to the Consciousness Discussion
% FFGFT Document 207 — Body
% Status: 8 May 2026
% ============================================================

\section{Introduction}\label{sec:intro}

The consciousness discussion has been part of the FFGFT series for some time. Two documents form its main lines:

\begin{itemize}
\item \textbf{Document 100 (\emph{Consciousness as fractal incoherence}):} Connection to the no-go theorem of Adlam, McQueen and Waegell~\cite{Adlam2025}, which shows that agency cannot arise in a purely unitary quantum system (world-model failure, deliberation impossibility, action-selection collapse). FFGFT's response: consciousness arises not from perfect quantum coherence but from \emph{structured fractal incoherence} --- a geometrically rooted deviation from unitarity, parametrised by $\xi = 4/30{,}000$.
\item \textbf{Document 170 (\emph{Consciousness as a threshold ladder}):} Consciousness as level~4 of a five-stage hierarchy of discrete complexity thresholds, carried by the fractal correction $K_{\text{frak}}^{(n)} = (\xi/(1+\xi))^{n/2}$ as a threshold operator. Self-declared as a working hypothesis.
\end{itemize}

Document~206 (\emph{Triangle--Matrix Reduction Theorem}, also dated 8 May 2026), with its proof of $\det(A_\xi) = 1-\xi$ and the Z$_3^3$ scale ladder $(1-\xi)^{n/3}$, has established an algebraic structure that grounds the phenomenological statements of Documents~100 and~170 structurally. In parallel, the discussion in the \emph{Information Physics Institute} mailing list in the weeks before this document has coined two central concepts that translate directly into the same structure: Onur Tekermen's \emph{constitutive openness} and Phillips Paul's \emph{Information Grounding Law} (IGL) within the Self-Dual~C-25 construction.

\medskip

The task of this document is limited: we show that the language used to discuss consciousness in FFGFT and in the IPI discussion is consistent with the algebraic skeleton established in Doc~206. We formulate no new theory of consciousness. We open no new research questions that are not already in Doc~100 or Doc~170. We merely check what the algebra structurally supports.

\begin{important}[What this document is and is not]
\textbf{It is:} A connection paper showing where Doc~206 algebraically meets the phenomenological line of Doc~100 and Doc~170.

\textbf{It is not:} A derivation of consciousness from $\xi$. The algebra enforces structural features (discreteness, openness, self-inversion); it does not identify any level as ``conscious''.
\end{important}

\section{What Document 206 algebraically achieves (summary)}\label{sec:doc206_summary}

For the connection argument we summarise the results from Doc~206 relevant for this document.

\subsection{\texorpdfstring{Z$_3$ triangle and 3$\times$3 matrix}{Z3 triangle and 3x3 matrix}}

The Z$_3$ cycle has adjacency matrix
\[
A = \begin{pmatrix} 0 & 0 & 1 \\ 1 & 0 & 0 \\ 0 & 1 & 0 \end{pmatrix},
\quad
A^3 = I, \quad \det(A) = 1, \quad \mathrm{Tr}(A) = 0.
\]
The eigenvalues are the third roots of unity $\{1, \omega, \omega^2\}$.

\subsection{\texorpdfstring{$\xi$-perturbation and perturbed adjacency}{xi-perturbation and perturbed adjacency}}

The closing edge is damped by $\xi$:
\[
A_\xi = \begin{pmatrix} 0 & 0 & 1-\xi \\ 1 & 0 & 0 \\ 0 & 1 & 0 \end{pmatrix},
\quad
A_\xi^3 = (1-\xi)\,I, \quad \det(A_\xi) = 1-\xi.
\]
The fractal dimension follows in linear order: $D_f = 3 - \xi$.

\subsection{\texorpdfstring{Z$_3^3$ tensor product and scale ladder}{Z3-cubed tensor product and scale ladder}}

The tensor product $T(\xi) = A_\xi \otimes A_\xi \otimes A_\xi$ has $27$ eigenvalues of modulus $(1-\xi)$ and produces the scale ladder
\[
\lambda_n = (1-\xi)^{n/3}, \quad n = 1, 2, 3, \dots
\]
These third-power exponents are not arbitrary but a consequence of the threefold tensor structure.

\subsection{Six bridges to other formulations}

Doc~206 tests six bridges (Austin, Tekermen, Phillips, Porter, Chekanov-Hakan, Moseley). Relevant to the consciousness question are \textbf{Tekermen} (constitutive openness) and \textbf{Phillips} (IGL as self-dual condition). The other bridges are not pursued here.

\section{Bridge A: Constitutive openness and world-relation}\label{sec:bridgeA}

\subsection{Tekermen's postulate in the IPI discussion}

Onur Tekermen, in his Unified Information Field framework (UIFT), formulates the thesis that a recursion is \emph{informationally meaningful} only if it does not close. A fully closed system has no world, no outside, no reference. Tekermen calls this \emph{constitutive openness}.

\subsection{Algebraic identification in FFGFT}

In FFGFT this corresponds exactly to the Non-closure Theorem from Doc~193, anchored algebraically in Doc~206 \S6:
\begin{equation}
\det(A_\xi) = 1 - \xi.
\end{equation}
For $\xi = 0$ we would have $\det = 1$, the system would be exactly cyclic (perfectly unitary in the matrix reading). Only $\xi > 0$ algebraically opens the cycle.

\begin{insightBox}[Algebraic form of Tekermen's postulate]
Constitutive openness $\equiv$ $\xi > 0$. The recursion is meaningful if and only if $\det(A_\xi) < 1$.
\end{insightBox}

\subsection{Connection to Document 100 and to Adlam's no-go theorem}

Adlam, McQueen and Waegell show that a purely unitary quantum system cannot carry agency. Document~100 responds: consciousness arises not from quantum coherence, but from fractal incoherence.

In the language of Doc~206 this is algebraically readable:

\begin{itemize}
\item \textbf{Adlam's no-go:} unitary system $=$ closed recursion $=$ $\det = 1$, hence $\xi = 0$.
\item \textbf{FFGFT response:} $\xi > 0$ geometrically forces a world-relation. The system has an outside because its algebraic skeleton does not close.
\end{itemize}

Adlam's no-go theorem and FFGFT's non-closure theorem are not two theorems but two sides of the same structural statement. Adlam shows what does \emph{not} work in a closed system; FFGFT shows what is the case in an open system.

\begin{interpretation}
This identification does not solve the consciousness problem. It does not say that $\xi > 0$ \emph{produces} consciousness. It says only that the structural precondition for world-reference, demanded by Adlam and postulated by Tekermen, is algebraically satisfied in FFGFT. What is built on top of this precondition remains open.
\end{interpretation}

\section{Bridge B: Self-dual and recursive feedback}\label{sec:bridgeB}

\subsection{Phillips' Information Grounding Law}

Phillips Paul postulates an algebraic condition for \emph{informationally anchored} structures within the Self-Dual~C-25 construction: an operation $A$ is \emph{self-dual} if
\begin{equation}
A^2 = A^\top = A^{-1}.
\end{equation}
Phillips calls this the \emph{Information Grounding Law}: information is anchored only if the underlying operation can refer back to itself without information loss.

\subsection{Algebraic identification in FFGFT}

In Doc~206 \S7 it is shown that
\begin{equation}
A^2 = A^\top \quad \text{for the Z}_3 \text{ adjacency matrix},
\end{equation}
and more generally $A^k = A^{\top, k}$ for all integer $k$. The Z$_3$ self-inversion is the elementary algebraic embodiment of Phillips' IGL.

\subsection{Connection to Document 170: recursive sensory feedback}

Document~170 names \emph{recursive sensory feedback} as the first criterion for the complexity threshold of level~2 (cellular self-maintenance): membrane receptors, chemical gradients, and intracellular signal cascades form a closed feedback loop. The cell does not passively perceive its environment but modulates its own response on the basis of past states.

In the language of Doc~206 this is the operative meaning of $A^2 = A^\top$: the operation ``perceive'' is invertible to itself --- perceiving and reacting are not two separate processes but the same Z$_3$ operation in different directions.

\begin{insightBox}[Algebraic form of recursive feedback]
Recursive sensory feedback $\equiv$ self-dual ($A^2 = A^\top$). Phillips' IGL and Doc~170's recursive sensory feedback section describe the same algebraic property.
\end{insightBox}

\subsection{What the bridge achieves --- and what it does not}

It achieves: it shows that the phenomenology in Doc~170 and the postulate from the IPI discussion sit on the same Z$_3$ skeleton. ``Recursive sensory feedback'' is not poetic or metaphorical --- it has a concrete algebraic form.

It does \emph{not} achieve: it does not say that all Z$_3$ self-inversions carry consciousness. A rotating screw satisfies $A^2 = A^\top$ and is not conscious. Self-duality is a \emph{necessary} structural condition (per Doc~170 and Phillips' IGL); it is not sufficient.

\section{\texorpdfstring{Bridge C: Z$_3^3$ scale ladder and discrete thresholds}{Bridge C: Z3-cubed scale ladder and discrete thresholds}}\label{sec:bridgeC}

\subsection{Doc 170: table of complexity levels}

Document~170 proposes a five-stage hierarchy:

\begin{center}
\begin{tabular}{c|l|l}
\textbf{Level} & \textbf{System} & \textbf{Threshold characteristic} \\
\hline
0 & Elementary particles & $T = 1/m$ as passive time field \\
1 & Atoms / crystals & Discrete geometric order \\
2 & Cell & Recursive self-maintenance, metabolism \\
3 & Nervous system & Integrated sensory feedback \\
4 & Consciousness & Self-modelling, world-reference \\
\end{tabular}
\end{center}

Doc~170 explicitly marks this as a working hypothesis.

\subsection{Doc 206: algebraic scale ladder}

Doc~206 \S4 shows that the Z$_3^3$ tensor structure has eigenvalues
\[
\lambda_n = (1-\xi)^{n/3}, \quad n \in \mathbb{Z}_{\geq 0}.
\]
These third-power exponents are not chosen arbitrarily but algebraically forced by the threefold tensor structure.

\subsection{What the algebra contributes to Doc 170}

\begin{itemize}
\item \textbf{Discreteness:} the levels in Doc~170 are discrete. This is consistent with the Z$_3^3$ scale ladder: $(1-\xi)^{n/3}$ is algebraically anchored only for $n = 0, 1, 2, 3, \dots$. Continuous intermediate values have no geometric status.
\item \textbf{Hierarchy:} larger $n$ corresponds to smaller modulus $(1-\xi)^{n/3}$, hence deeper recursion. This is compatible with the picture that more complex structures are more deeply embedded.
\item \textbf{Topological forcing:} Doc~170 says the levels are ``topologically forced''. This is algebraically anchored: the $\omega^k$ phases with $k = 0, 1, 2$ are not arbitrary but exactly the Z$_3$ characters.
\end{itemize}

\subsection{What the algebra does not contribute to Doc 170}

\begin{important}[Limit of connection]
The Z$_3^3$ scale ladder enforces \emph{discreteness} but not the concrete \emph{level assignment} from Doc~170. The algebra does not say ``level~4 is consciousness''. Which $n$ corresponds to which biological level remains a phenomenological question and is marked as a working hypothesis in Doc~170.
\end{important}

The third powers generate infinitely many layers, not just five. Doc~170 chooses five because the biological phenomenology suggests these levels, not because the algebra exhibits them.

\section{Methodological self-check}\label{sec:methodology}

In Doc~206 two self-check sections were introduced: a numerology filter (\S12) and a methodological status section on cosmological data (\S11). The same discipline applies to the consciousness question.

\subsection{Numerology filter}

We check whether our bridges have the status of a theoretical finding or whether they could appear coincidental.

\begin{itemize}
\item \textbf{Bridge A (Tekermen $\equiv$ $\xi > 0$):} algebraically identical, not numerological. Status: theoretical.
\item \textbf{Bridge B (Phillips IGL $\equiv$ $A^2 = A^\top$):} algebraically identical (Doc~206 \S7). Status: theoretical.
\item \textbf{Bridge C (Doc 170 levels $\equiv$ $(1-\xi)^{n/3}$):} \emph{Only partially} theoretical. Discreteness and hierarchy are algebraically anchored. The concrete level assignment is \emph{not} algebraically justified and remains phenomenological.
\end{itemize}

\subsection{What we do not claim}

We avoid four typical fallacies:

\begin{enumerate}
\item \textbf{Numerological consciousness identification.} Statements such as ``consciousness corresponds to layer $n = 8$ because $(1-\xi)^{8/3}$ has a particular order of magnitude'' are not admissible. There is no geometric indication of a privileged $n$ value.
\item \textbf{Circular Adlam resolution.} We do not claim that $\xi > 0$ ``solves'' the Adlam problem. We say only that it satisfies the structural precondition that Adlam misses. What stands above this precondition is open.
\item \textbf{IGL-to-consciousness leap.} Consciousness does not follow from $A^2 = A^\top$. Self-duality is necessary, not sufficient.
\item \textbf{Level fixation.} Doc~170 assigns five levels. Doc~206 generates infinitely many. We leave open which levels are ``occupied'' and whether the Doc~170 assignment is exactly right.
\end{enumerate}

\subsection{Methodological status of the consciousness question}

Doc~206 is a proven theorem with reproducible scripts. Doc~100 and Doc~170 are explicitly formulated as connection papers and working hypotheses. Doc~207 (this document) is a synthesis paper; it provides no new proofs but establishes structural identifications.

\begin{interpretation}
The consciousness question is not closed within the FFGFT series. Doc~206 makes the language of the discussion algebraically consistent, but the ontological question ``what is consciousness?'' remains open. Anyone attempting to derive a theory of consciousness from FFGFT must clearly separate structural statements (algebraically supported) from phenomenological statements (working hypothesis).
\end{interpretation}

\section{Empirical attachability: predictions under explicit auxiliary assumptions}\label{sec:predictions}

The algebra in Doc~206 is a pure structural statement; it makes no direct predictions about measurable consciousness markers. If, however, we take the \emph{discreteness hypothesis} from Document~170 (consciousness levels are topologically forced, not gradually reachable) as an auxiliary assumption, four qualitative empirical predictions follow. We formulate them in the form Doc~206 \S12 permits: \emph{structural features without numerical values}.

\subsection{What follows from FFGFT and what does not}\label{sec:predictions_follows}

\begin{important}[Separating structural statements from numerical values]
What \emph{follows} from FFGFT: if the discreteness hypothesis holds, consciousness markers should show qualitatively discontinuous structures (bimodality, jumps, tipping points, taxonomic thresholds).

What \emph{does not follow} from FFGFT: concrete numerical values for cutoffs, critical doses, threshold ages, or which species cross which threshold. These depend on neurobiological phenomenology, not on $\xi$.
\end{important}

\subsection{Prediction A: bimodality of neural consciousness markers}\label{sec:predictionA}

If the discreteness hypothesis holds, a consciousness marker such as the Perturbational Complexity Index (PCI) should not show a continuous distribution but a \emph{bimodal} separation into two clusters with a gap between them.

\begin{itemize}
\item \textbf{Follows:} bimodality as a qualitative structural feature.
\item \textbf{Does not follow:} a concrete cutoff (e.\,g.\ PCI~$\approx 0.31$ as used clinically) or a gap size in $\xi$ units. Statements of the form ``the gap is proportional to $\xi$'' have no algebraic backing.
\end{itemize}

\textbf{Empirical status:} Casali et al.~(2013)~\cite{Casali2013} have established a clinical PCI cutoff. Whether the distribution is in fact bimodal (real gap) or unimodal with a pragmatic cutoff is an open empirical question.

\subsection{Prediction B: developmental jump rather than gradient}\label{sec:predictionB}

Consciousness development in infants should show a jump, not a continuous gradient.

\begin{itemize}
\item \textbf{Follows:} jump character as a structural feature.
\item \textbf{Does not follow:} the timing of the jump. This depends on neural maturation (e.\,g.\ myelination of recurrent pathways), not on $\xi$.
\end{itemize}

\subsection{Prediction C: anaesthesia tipping rather than gradation}\label{sec:predictionC}

Under slowly increasing anaesthetic dose, loss of consciousness should occur at a tipping point, not gradually.

\begin{itemize}
\item \textbf{Follows:} tipping character.
\item \textbf{Does not follow:} the critical dose, substance specificity, dose-response time constants.
\end{itemize}

\subsection{Prediction D: species classification rather than scaling}\label{sec:predictionD}

Consciousness should not be a function of complexity alone. Increased complexity (more neurons, more connections) without structural Z$_3$ self-inversion should not cross any consciousness threshold.

\begin{itemize}
\item \textbf{Follows:} the existence of a structural threshold that cannot be overcome by scaling alone.
\item \textbf{Does not follow:} which taxa lie above or below the threshold. This is to be answered neurobiologically; it is not derivable from $\xi$.
\end{itemize}

\textbf{Consequence for AI:} pure scaling of transformer architectures does not, per Prediction~D, reach consciousness. What would be necessary is a Z$_3$-equivalent self-inversion with metabolic embodiment in the sense of Doc~170 \S2 (vulnerability, metabolism, $T = 1/m$). This is a structural statement, not a temporal forecast about future AI development.

\subsection{What we do not predict (filter analogous to Doc 206 \S12)}\label{sec:predictions_filter}

We delimit explicitly what \emph{cannot} be derived from FFGFT, even if it appears tempting:

\begin{enumerate}
\item \textbf{No quantitative PCI threshold from $\xi$.} Statements of the form ``$C_{\text{crit}} = \kappa \cdot \xi$ with $\kappa = \mathcal{O}(1)$'' are numerology. A free factor $\kappa$ permits any desired value.
\item \textbf{No derivation of the EEG spectral exponent.} Proposals such as $\beta = 2/3 + \xi/(2\pi)$ contain factors ($1/(2\pi)$) that are not justified by the FFGFT algebra. The effect would be smaller than EEG noise and thus practically unfalsifiable.
\item \textbf{No identity thesis.} Statements of the form ``consciousness $\equiv$ persistent recursive coupling with $T = 1/m$ modulation'' do not solve the Hard Problem; they redefine it. Self-duality, non-closure, and the $n/3$ hierarchy are necessary structural conditions --- not sufficient.
\item \textbf{No scale ladder $(1-\xi)^{k-1}$.} The only scale ladder algebraically supported by Doc~206 is $(1-\xi)^{n/3}$. It does not distinguish ``consciousness levels'' from ``non-consciousness levels''.
\item \textbf{No Z$_3$ identification with three anatomical brain areas.} Doc~206 has a Z$_3$ triangle with three nodes in an abstract scale space. Identification with concrete anatomy (V1 $\to$ V4 $\to$ prefrontal) is a category shift without algebraic backing.
\end{enumerate}

\subsection{Falsifiability}\label{sec:predictions_falsification}

The four qualitative predictions (A--D) are falsifiable:

\begin{itemize}
\item If PCI distributions are empirically clearly continuous, Prediction~A fails.
\item If infant consciousness markers show a cleanly continuous gradient, Prediction~B fails.
\item If anaesthesia effects are exactly gradual, Prediction~C fails.
\item If complexity alone (e.\,g.\ in scaled AI) produces consciousness, Prediction~D fails.
\end{itemize}

\begin{interpretation}
Falsification of these qualitative predictions falsifies the \emph{discreteness hypothesis from Doc~170}, not the algebraic skeleton from Doc~206. The Z$_3$ algebra is correct or incorrect independently of the consciousness question.
\end{interpretation}

\section{Relation to the IPI mailing list discussion}\label{sec:ipi}

In the weeks before this document, several central concepts were proposed in the IPI mailing list that now fit directly into the Z$_3$ skeleton from Doc~206.

\begin{itemize}
\item \textbf{Tekermen's ``open fraction''} (May 2026): the idea that consciousness is a \emph{non-binary spectrum} corresponds to the Z$_3^3$ scale ladder. ``Non-binary'' algebraically means: not only $0$ and $1$ but $(1-\xi)^{n/3}$ for all $n$.
\item \textbf{Phillips' IGL and Self-Dual~C-25} (May 2026): algebraically identified with $A^2 = A^\top$ in Doc~206 \S7.
\item \textbf{Austin's $D = 3 + \varepsilon$} (May 2026): identified in Doc~206 \S5 as the mirror image of $D_f = 3 - \xi$ ($\varepsilon = -\xi$). This has no direct bearing on the consciousness question, but it is methodologically relevant: the same geometric structure can be formulated with different sign conventions.
\end{itemize}

\subsection{What the IPI discussion and Doc 207 achieve}

The IPI mailing list discussion in the weeks before this document has reached the language level at which algebraic identifications become possible. Doc~207 is not the beginning of this discussion but a consolidation.

The language now used to discuss consciousness (open recursion, self-dual, non-binary spectrum) is consistent with the Z$_3$ geometry. This is methodological progress --- not content progress.

\section{Balance}\label{sec:balance}

\begin{conclusionBox}[Balance: what Doc 207 establishes]
Three structural bridges between Doc~206 and the existing consciousness discussion are algebraically supported:

\begin{enumerate}
\item \textbf{Constitutive openness} (Tekermen) $\equiv$ $\xi > 0$ (Non-closure Theorem from Doc~193, anchored algebraically in Doc~206 \S6).
\item \textbf{Information Grounding Law} (Phillips) $\equiv$ $A^2 = A^\top$ (Z$_3$ self-inversion in Doc~206 \S7).
\item \textbf{Discrete complexity levels} (Doc~170) algebraically consistent with the Z$_3^3$ scale ladder $(1-\xi)^{n/3}$ --- discreteness enforced, level assignment phenomenological.
\end{enumerate}

Adlam's no-go theorem and FFGFT's non-closure theorem are two sides of the same structural statement.
\end{conclusionBox}

\subsection{What remains open}

\begin{itemize}
\item The ontological consciousness question (the ``Hard Problem'') is not touched by any of the three bridges.
\item The concrete assignment of Doc~170 levels to $n$ values is not algebraically justified.
\item Sufficient conditions for consciousness remain open; Doc~206 supplies only necessary structural conditions for the language of the discussion.
\item The cell as ``first level of consciousness'' (Doc~170) is a working hypothesis, not a result of Doc~206.
\end{itemize}

\subsection{Relation within the corpus}

\begin{center}
\begin{tabular}{c|p{8cm}|l}
\textbf{Document} & \textbf{Contribution to consciousness question} & \textbf{Status} \\
\hline
100 & Consciousness as fractal incoherence, Adlam connection & Connection paper \\
170 & Threshold ladder, $K_{\text{frak}}$ as threshold operator & Working hypothesis \\
193 & Non-closure theorem (anchored algebraically in 206) & Proven \\
206 & Z$_3$ skeleton, six bridges, $D_f = 3 - \xi$ & Proven \\
\textbf{207} & \textbf{Structural bridges between 206 and 100/170} & \textbf{Synthesis} \\
\end{tabular}
\end{center}

Doc~207 makes no new statement about consciousness. It shows that the existing statements from Doc~100, Doc~170 and the IPI discussion sit within a consistent algebraic framework. The consciousness question itself remains an open research question. Doc~206 supplies the language in which further work can be precisely formulated; it does not supply an answer to what consciousness is.

\begin{philosophical}[Closing]
The question of what consciousness is is not answered by an algebraic bridge. But the question of in which language one can ask it without falling into numerological or merely metaphorical dead ends has, through Doc~206, gained sharper form. Nothing more is gained. Nothing less is intended.
\end{philosophical}

\section*{Acknowledgements}

The bridges formulated here build on the discussions in the \emph{Information Physics Institute} mailing list and on the work of Onur Tekermen, Phillips Paul, and Peter Austin. The phenomenological line owes itself to engagement with the no-go theorem of Adlam, McQueen and Waegell~\cite{Adlam2025}.

\begin{thebibliography}{9}
\bibitem{Adlam2025} E.\,C.~Adlam, K.\,J.~McQueen, M.~Waegell. \emph{Agency cannot be a purely quantum phenomenon.} arXiv:2510.13247 (2025).
\bibitem{Casali2013} A.\,G.~Casali, O.~Gosseries, M.~Rosanova, M.~Boly, S.~Sarasso, K.\,R.~Casali, S.~Casarotto, M.-A.~Bruno, S.~Laureys, G.~Tononi, M.~Massimini. \emph{A theoretically based index of consciousness independent of sensory processing and behavior.} \emph{Sci.\ Transl.\ Med.}\ \textbf{5}, 198ra105 (2013).
\end{thebibliography}

\section*{License and availability}

\textbf{GitHub:} \href{https://github.com/jpascher/T0-Time-Mass-Duality}{\texttt{github.com/jpascher/T0-Time-Mass-Duality}}\\
\textbf{Zenodo (predecessor v1.0.13):} \href{https://zenodo.org/records/20041529}{\texttt{zenodo.org/records/20041529}} (DOI: \href{https://doi.org/10.5281/zenodo.20041529}{10.5281/zenodo.20041529})\\
\textbf{ORCID:} \href{https://orcid.org/0009-0000-6518-4064}{0009-0000-6518-4064}\\
\textbf{License:} CC-BY-SA 4.0


\end{document}
